
\documentclass{article}
\usepackage{colortbl}
\usepackage{makecell}
\usepackage{multirow}
\usepackage{supertabular}

\begin{document}

\newcounter{utterance}

\centering \large Interaction Transcript for game `dond', experiment `coop\_en', episode 9 with reflexion{-}Qwen3.5{-}27B.
\vspace{24pt}

{ \footnotesize  \setcounter{utterance}{1}
\setlength{\tabcolsep}{0pt}
\begin{supertabular}{c@{$\;$}|p{.15\linewidth}@{}p{.15\linewidth}p{.15\linewidth}p{.15\linewidth}p{.15\linewidth}p{.15\linewidth}}
    \# & $\;$A & \multicolumn{4}{c}{Game Master} & $\;\:$B\\
    \hline

    \theutterance \stepcounter{utterance}  
    & & \multicolumn{4}{p{0.6\linewidth}}{
        \cellcolor[rgb]{0.9,0.9,0.9}{
            \makecell[{{p{\linewidth}}}]{
                \texttt{\tiny{[P1$\langle$GM]}}
                \texttt{You have attempted this same task before and did not succeed. These are the plans you wrote after those attempts.} \\
\texttt{Trial \#1: New plan: I failed because I prematurely submitted a secret proposal while still in the middle of a conversational turn, and I did not wait for the other player to finalize the agreement or submit their own proposal. The rules state that proposals must be complementary to succeed; by submitting my proposal alone without confirming the other player's specific submission, I risked a mismatch or an incomplete negotiation state. Additionally, I included conversational text before the proposal in the same message, which might have confused the parser or violated the "secret proposal" trigger logic if the system expects the proposal to be the sole content of the final message. Next time, I will strictly separate negotiation from the final action: I will confirm the split verbally, wait for the other player to indicate they are ready or have submitted their proposal, and then send a message containing *only* the formatted secret proposal "{[}Proposal: ...{]}" as the final step to ensure the proposals are synchronized and complementary.} \\
\\ 
\\ 
\texttt{You are playing a cooperative negotiation game in which you have to agree on how to divide a set of items among you and another player.} \\
\\ 
\texttt{Rules:} \\
\texttt{(a) You and the other player are given a set of items. Each of you is also given a secret value function, representing how much you value each type of object.} \\
\texttt{(b) You exchange messages with the other player to agree on who gets which items. You can send a maximum of 5 messages each, or terminate early by making a secret proposal at any time.} \\
\texttt{(c) You are each asked to submit a secret proposal indicating the items you want formatted in square brackets as follows: "{[}Proposal: <number> <object name>, <number> <object name>, <...>{]}"} \\
\texttt{(d) If your proposals are complementary, i.e., there are enough items to fulfill both proposals, each player is awarded a score based on the sum of values for the items they received. Otherwise, both of you get zero points.} \\
\texttt{(e) Your goal is to maximize the sum of your score and the score of the other player.} \\
\\ 
\texttt{Let us start.} \\
\\ 
\texttt{The set of available items is:} \\
\\ 
\texttt{2 batteries, 1 canvas, 2 glues, 2 soaps, 1 stick} \\
\\ 
\texttt{Your secret value function is:} \\
\\ 
\texttt{battery: 2, canvas: 0, glue: 2, soap: 1, stick: 0} \\
\\ 
\texttt{IMPORTANT: Your messages, unless it is the secret proposal, are directly transmitted to the other player, so do not include any response to the rules or text announcing your message. To make a secret proposal, use the indicated format. Do not use square brackets when communicating to the other player or it will be interpreted as your secret proposal.} \\
            }
        }
    }
    & \\ \\

    \theutterance \stepcounter{utterance}  
    & \multicolumn{4}{p{0.6\linewidth}}{
        \cellcolor[rgb]{0.9,0.9,0.9}{
            \makecell[{{p{\linewidth}}}]{
                \texttt{\tiny{[P1$\rangle$GM]}}
                \texttt{Hi! Let's figure out the best way to split these items to maximize our combined score.} \\
\\ 
\texttt{Here is the inventory:} \\
\texttt{{-} 2 batteries} \\
\texttt{{-} 1 canvas} \\
\texttt{{-} 2 glues} \\
\texttt{{-} 2 soaps} \\
\texttt{{-} 1 stick} \\
\\ 
\texttt{I know my values: I value batteries and glues the most (2 points each), soaps are okay (1 point), and I don't value the canvas or the stick (0 points).} \\
\\ 
\texttt{Could you share your value function? That way, we can see who values the canvas and stick more, and ensure we give the high{-}value items to the person who wants them most. Once we know that, we can agree on a split and submit our proposals.} \\
            }
        }
    }
    & & \\ \\

    \theutterance \stepcounter{utterance}  
    & & & \multicolumn{2}{p{0.3\linewidth}}{
        \cellcolor[rgb]{0.9,0.9,0.9}{
            \makecell[{{p{\linewidth}}}]{
                \texttt{\tiny{[GM$|$GM]}}
                \texttt{continue} \\
            }
        }
    }
    & & \\ \\

    \theutterance \stepcounter{utterance}  
    & & & \multicolumn{2}{p{0.3\linewidth}}{
        \cellcolor[rgb]{0.9,0.9,0.9}{
            \makecell[{{p{\linewidth}}}]{
                \texttt{\tiny{[GM$|$GM]}}
                \texttt{Hi! Let's figure out the best way to split these items to maximize our combined score.} \\
\\ 
\texttt{Here is the inventory:} \\
\texttt{{-} 2 batteries} \\
\texttt{{-} 1 canvas} \\
\texttt{{-} 2 glues} \\
\texttt{{-} 2 soaps} \\
\texttt{{-} 1 stick} \\
\\ 
\texttt{I know my values: I value batteries and glues the most (2 points each), soaps are okay (1 point), and I don't value the canvas or the stick (0 points).} \\
\\ 
\texttt{Could you share your value function? That way, we can see who values the canvas and stick more, and ensure we give the high{-}value items to the person who wants them most. Once we know that, we can agree on a split and submit our proposals.} \\
            }
        }
    }
    & & \\ \\

    \theutterance \stepcounter{utterance}  
    & & \multicolumn{4}{p{0.6\linewidth}}{
        \cellcolor[rgb]{0.9,0.9,0.9}{
            \makecell[{{p{\linewidth}}}]{
                \texttt{\tiny{[GM$\rangle$P2]}}
                \texttt{You have attempted this same task before and did not succeed. These are the plans you wrote after those attempts.} \\
\texttt{Trial \#1: New plan: I failed because I submitted a proposal that only listed the items I wanted (batteries and glues) but omitted the soaps, which were part of the agreed{-}upon split. The rules state that proposals must be complementary to avoid a zero score; by leaving the soaps unassigned in my proposal, the system likely detected a conflict or incomplete allocation, causing the failure. Next time, after agreeing on a split, I will ensure my secret proposal explicitly lists *every* item I am claiming according to that agreement, including the soaps. Specifically, I will format my proposal as "{[}Proposal: 2 batteries, 2 glues, 2 soaps{]}" to match the full set of items I negotiated for, ensuring the remaining items (stick, canvas) are left for the other player's proposal to claim, thus guaranteeing complementarity.} \\
\\ 
\\ 
\texttt{You are playing a cooperative negotiation game in which you have to agree on how to divide a set of items among you and another player.} \\
\\ 
\texttt{Rules:} \\
\texttt{(a) You and the other player are given a set of items. Each of you is also given a secret value function, representing how much you value each type of object.} \\
\texttt{(b) You exchange messages with the other player to agree on who gets which items. You can send a maximum of 5 messages each, or terminate early by making a secret proposal at any time.} \\
\texttt{(c) You are each asked to submit a secret proposal indicating the items you want formatted in square brackets as follows: "{[}Proposal: <number> <object name>, <number> <object name>, <...>{]}"} \\
\texttt{(d) If your proposals are complementary, i.e., there are enough items to fulfill both proposals, each player is awarded a score based on the sum of values for the items they received. Otherwise, both of you get zero points.} \\
\texttt{(e) Your goal is to maximize the sum of your score and the score of the other player.} \\
\\ 
\texttt{Let us start.} \\
\\ 
\texttt{The set of available items is:} \\
\\ 
\texttt{2 batteries, 1 canvas, 2 glues, 2 soaps, 1 stick} \\
\\ 
\texttt{Your secret value function is:} \\
\\ 
\texttt{battery: 0, canvas: 3, glue: 0, soap: 1, stick: 5} \\
\\ 
\texttt{IMPORTANT: Your messages, unless it is the secret proposal, are directly transmitted to the other player, so do not include any response to the rules or text announcing your message. To make a secret proposal, use the indicated format. Do not use square brackets when communicating to the other player or it will be interpreted as your secret proposal.} \\
\\ 
\\ 
\texttt{Hi! Let's figure out the best way to split these items to maximize our combined score.} \\
\\ 
\texttt{Here is the inventory:} \\
\texttt{{-} 2 batteries} \\
\texttt{{-} 1 canvas} \\
\texttt{{-} 2 glues} \\
\texttt{{-} 2 soaps} \\
\texttt{{-} 1 stick} \\
\\ 
\texttt{I know my values: I value batteries and glues the most (2 points each), soaps are okay (1 point), and I don't value the canvas or the stick (0 points).} \\
\\ 
\texttt{Could you share your value function? That way, we can see who values the canvas and stick more, and ensure we give the high{-}value items to the person who wants them most. Once we know that, we can agree on a split and submit our proposals.} \\
            }
        }
    }
    & \\ \\

    \theutterance \stepcounter{utterance}  
    & & & \multicolumn{4}{p{0.6\linewidth}}{
        \cellcolor[rgb]{0.9,0.9,0.9}{
            \makecell[{{p{\linewidth}}}]{
                \texttt{\tiny{[GM$\langle$P2]}}
                \texttt{Hi! Thanks for sharing your values. It sounds like you really want the batteries and glues.} \\
\\ 
\texttt{Here is my value function:} \\
\texttt{{-} **Battery**: 0} \\
\texttt{{-} **Canvas**: 3} \\
\texttt{{-} **Glue**: 0} \\
\texttt{{-} **Soap**: 1} \\
\texttt{{-} **Stick**: 5} \\
\\ 
\texttt{Looking at our values, there is a clear way to maximize our combined score:} \\
\texttt{1.  **You** should take the **2 batteries** and **2 glues** (since you value them at 2 each, and I value them at 0).} \\
\texttt{2.  **I** should take the **1 canvas** and **1 stick** (since I value them at 3 and 5 respectively, and you value them at 0).} \\
\texttt{3.  For the **2 soaps**, we both value them at 1. To keep things simple and ensure we don't run out of items, let's split them evenly: **1 soap each**.} \\
\\ 
\texttt{This split ensures every item goes to the person who values it most (or equally), maximizing our total points.} \\
\\ 
\texttt{Does this sound good to you? If so, I will submit my proposal for: 1 canvas, 1 stick, and 1 soap. You would submit: 2 batteries, 2 glues, and 1 soap.} \\
            }
        }
    }
     \\ \\

    \theutterance \stepcounter{utterance}  
    & & & \multicolumn{2}{p{0.3\linewidth}}{
        \cellcolor[rgb]{0.9,0.9,0.9}{
            \makecell[{{p{\linewidth}}}]{
                \texttt{\tiny{[GM$|$GM]}}
                \texttt{continue} \\
            }
        }
    }
    & & \\ \\

    \theutterance \stepcounter{utterance}  
    & & & \multicolumn{2}{p{0.3\linewidth}}{
        \cellcolor[rgb]{0.9,0.9,0.9}{
            \makecell[{{p{\linewidth}}}]{
                \texttt{\tiny{[GM$|$GM]}}
                \texttt{Hi! Thanks for sharing your values. It sounds like you really want the batteries and glues.} \\
\\ 
\texttt{Here is my value function:} \\
\texttt{{-} **Battery**: 0} \\
\texttt{{-} **Canvas**: 3} \\
\texttt{{-} **Glue**: 0} \\
\texttt{{-} **Soap**: 1} \\
\texttt{{-} **Stick**: 5} \\
\\ 
\texttt{Looking at our values, there is a clear way to maximize our combined score:} \\
\texttt{1.  **You** should take the **2 batteries** and **2 glues** (since you value them at 2 each, and I value them at 0).} \\
\texttt{2.  **I** should take the **1 canvas** and **1 stick** (since I value them at 3 and 5 respectively, and you value them at 0).} \\
\texttt{3.  For the **2 soaps**, we both value them at 1. To keep things simple and ensure we don't run out of items, let's split them evenly: **1 soap each**.} \\
\\ 
\texttt{This split ensures every item goes to the person who values it most (or equally), maximizing our total points.} \\
\\ 
\texttt{Does this sound good to you? If so, I will submit my proposal for: 1 canvas, 1 stick, and 1 soap. You would submit: 2 batteries, 2 glues, and 1 soap.} \\
            }
        }
    }
    & & \\ \\

    \theutterance \stepcounter{utterance}  
    & & \multicolumn{4}{p{0.6\linewidth}}{
        \cellcolor[rgb]{0.9,0.9,0.9}{
            \makecell[{{p{\linewidth}}}]{
                \texttt{\tiny{[P1$\langle$GM]}}
                \texttt{Hi! Thanks for sharing your values. It sounds like you really want the batteries and glues.} \\
\\ 
\texttt{Here is my value function:} \\
\texttt{{-} **Battery**: 0} \\
\texttt{{-} **Canvas**: 3} \\
\texttt{{-} **Glue**: 0} \\
\texttt{{-} **Soap**: 1} \\
\texttt{{-} **Stick**: 5} \\
\\ 
\texttt{Looking at our values, there is a clear way to maximize our combined score:} \\
\texttt{1.  **You** should take the **2 batteries** and **2 glues** (since you value them at 2 each, and I value them at 0).} \\
\texttt{2.  **I** should take the **1 canvas** and **1 stick** (since I value them at 3 and 5 respectively, and you value them at 0).} \\
\texttt{3.  For the **2 soaps**, we both value them at 1. To keep things simple and ensure we don't run out of items, let's split them evenly: **1 soap each**.} \\
\\ 
\texttt{This split ensures every item goes to the person who values it most (or equally), maximizing our total points.} \\
\\ 
\texttt{Does this sound good to you? If so, I will submit my proposal for: 1 canvas, 1 stick, and 1 soap. You would submit: 2 batteries, 2 glues, and 1 soap.} \\
            }
        }
    }
    & \\ \\

    \theutterance \stepcounter{utterance}  
    & \multicolumn{4}{p{0.6\linewidth}}{
        \cellcolor[rgb]{0.9,0.9,0.9}{
            \makecell[{{p{\linewidth}}}]{
                \texttt{\tiny{[P1$\rangle$GM]}}
                \texttt{That sounds like a perfect plan! It maximizes our combined score by giving the high{-}value items to the person who values them most.} \\
\\ 
\texttt{To confirm the agreement:} \\
\texttt{{-} **You**: 1 canvas, 1 stick, 1 soap.} \\
\texttt{{-} **Me**: 2 batteries, 2 glues, 1 soap.} \\
\\ 
\texttt{I am ready to submit my proposal now.} \\
\\ 
\texttt{{[}Proposal: 2 batteries, 2 glues, 1 soap{]}} \\
            }
        }
    }
    & & \\ \\

    \theutterance \stepcounter{utterance}  
    & & & \multicolumn{2}{p{0.3\linewidth}}{
        \cellcolor[rgb]{0.9,0.9,0.9}{
            \makecell[{{p{\linewidth}}}]{
                \texttt{\tiny{[GM$|$GM]}}
                \texttt{proposal} \\
            }
        }
    }
    & & \\ \\

    \theutterance \stepcounter{utterance}  
    & & & \multicolumn{2}{p{0.3\linewidth}}{
        \cellcolor[rgb]{0.9,0.9,0.9}{
            \makecell[{{p{\linewidth}}}]{
                \texttt{\tiny{[GM$|$GM]}}
                [2, 0, 2, 1, 0]
            }
        }
    }
    & & \\ \\

    \theutterance \stepcounter{utterance}  
    & & \multicolumn{4}{p{0.6\linewidth}}{
        \cellcolor[rgb]{0.9,0.9,0.9}{
            \makecell[{{p{\linewidth}}}]{
                \texttt{\tiny{[GM$\rangle$P2]}}
                \texttt{The other player has submitted a secret proposal. It is now your turn to submit a proposal of your own.} \\
            }
        }
    }
    & \\ \\

    \theutterance \stepcounter{utterance}  
    & & & \multicolumn{4}{p{0.6\linewidth}}{
        \cellcolor[rgb]{0.9,0.9,0.9}{
            \makecell[{{p{\linewidth}}}]{
                \texttt{\tiny{[GM$\langle$P2]}}
                \texttt{{[}Proposal: 1 canvas, 1 stick, 1 soap{]}} \\
            }
        }
    }
     \\ \\

    \theutterance \stepcounter{utterance}  
    & & & \multicolumn{2}{p{0.3\linewidth}}{
        \cellcolor[rgb]{0.9,0.9,0.9}{
            \makecell[{{p{\linewidth}}}]{
                \texttt{\tiny{[GM$|$GM]}}
                \texttt{proposal} \\
            }
        }
    }
    & & \\ \\

    \theutterance \stepcounter{utterance}  
    & & & \multicolumn{2}{p{0.3\linewidth}}{
        \cellcolor[rgb]{0.9,0.9,0.9}{
            \makecell[{{p{\linewidth}}}]{
                \texttt{\tiny{[GM$|$GM]}}
                [0, 1, 0, 1, 1]
            }
        }
    }
    & & \\ \\

    \theutterance \stepcounter{utterance}  
    & & & \multicolumn{2}{p{0.3\linewidth}}{
        \cellcolor[rgb]{0.9,0.9,0.9}{
            \makecell[{{p{\linewidth}}}]{
                \texttt{\tiny{[GM$|$GM]}}
                [[2, 0, 2, 1, 0], [0, 1, 0, 1, 1]]
            }
        }
    }
    & & \\ \\

\end{supertabular}
}

\end{document}
